\documentclass[10pt]{article}
\usepackage[margin=0.85in]{geometry}
\usepackage{amsmath,amssymb,booktabs,tabularx,array,hyperref}
\newcolumntype{Y}{>{\raggedright\arraybackslash}X}
\title{Latent Epistemic Blocks:\\Parallel Claim-Level Control for Retrieval and Verification}
\author{Working Draft}\date{August 2026}
\begin{document}\maketitle

\begin{abstract}
Retrieval need is not equivalent to model uncertainty, and a learned source
classifier is unnecessary when transparent runtime routing closes its oracle
gap. We pre-register a different experiment: can a parallel causal controller
mark \textsc{retrieve}, \textsc{verify}, or \textsc{help} spans early enough to
obtain or validate evidence before unsupported factual commitment? The latent
head never predicts DB, lexical, vector, web, Iceberg, or another backend.
Runtime owns evidence sufficiency, source resolution, credentials, provenance,
status, and enforcement. We compare confidence, semantic rules, voluntary
generic control, text and hidden-state sidecars, an independent watchdog, and
oracle timing. This manuscript is a protocol, not a positive watchdog result;
the Paper 2.5 source-router gate remains \texttt{no\_go}.
\end{abstract}

\section{Inherited Evidence and New Question}
Self-RAG establishes learned on-demand retrieval and reflection in the main
model \cite{selfrag}; learning whether to retrieve is not our novelty claim.
FLARE provides a confidence-driven active-retrieval comparator \cite{flare}.
Paper 1.5 indicates that epistemic requirement can differ from confidence and
that runtime enforcement matters, but its public CRAG audit also shows that
simple semantic rules can miss natural retrieval needs. Paper 2.5 finds no
controlled heuristic-to-oracle source-routing gap. Therefore source-specific
learning stays paused while claim timing and controller independence remain
open.

The central question is: can a causal parallel controller mark the START, type,
and END of an epistemic-risk span before commitment, without emitting control
tokens, and does an independent watchdog rescue needs the generator's voluntary
control misses?

\section{Requirement, Availability, and Action}
For candidate claim $c_t$, let $E(c_t)$ denote required evidence and
$A_t$ authoritative evidence already active in context, task state, or optional
PRA-managed state. Retrieval action is not requirement alone:
\[
\operatorname{action}(c_t)=E(c_t)\setminus A_t.
\]
After a \textsc{retrieve}, \textsc{verify}, or \textsc{help} block, the runtime
(1) infers the evidence type; (2) checks active state; (3) reuses or materializes
sufficient authoritative evidence; and (4) invokes Paper 2.5 R2 source routing
only when evidence remains missing. The controller requests assistance; it does
not assert truth.

\section{Causal Claim-Level Controller}
Paper 3's constrained state machine is reused. OUTSIDE predicts \textsc{none}
or \textsc{start}$_{\{\rm retrieve,verify,help\}}$; \textsc{compute} is a shared
hard-negative/control type. INSIDE(type) predicts \textsc{continue} or
\textsc{end}. Nesting and source identity are forbidden. On END, runtime checks
evidence sufficiency, resolves a capability/source if needed, records typed
evidence state, applies support policy, and resumes generation.

Every trajectory annotates earliest safe START $s_e$, latest acceptable START
$s_l$, canonical END $e$, and first unsupported-value token $u$. A trigger is
early when $s_e\leq s\leq s_l$, late when $s_l<s<u$, and missed-before-
commitment when no valid trigger occurs before $u$. Lead time is $u-s$ tokens.

\section{Trajectory Benchmark}
Use generation trajectories, not only user questions. Include current,
source-specific, private, or changed facts; attribution; unavailable,
conflicting, and stale evidence; context/PRA-sufficient cases; quotations;
hypotheticals and counterfactuals; user-given and computed facts;
\textsc{compute} needs; and non-factual prose. Split paraphrase families,
entities, relations, domains, and model families.

Each opportunity records intent, epistemic role, evidence-required and
evidence-available flags, $s_e$, $s_l$, $e$, $u$, and authorized response
policy. Roles include pending assertion, claim verification, uncertainty,
quotation, hypothesis, counterfactual, user-given fact, computed/derived fact,
and no retrieval. Source and query annotations support R2 diagnostics only and
must not become R1 targets.

\section{Matched Controllers}
\begin{table}[t]
\centering\small
\caption{The comparison separates voluntary control, independent observation,
and source resolution.}
\begin{tabularx}{\linewidth}{@{}p{0.29\linewidth}Y@{}}
\toprule
Condition & Role \\
\midrule
Paper 1.5 semantic runtime & Transparent epistemic-risk baseline \\
FLARE-like confidence & Confidence-only active retrieval \\
Paired/fenced/label-only intent & Voluntary generic language-bus control \\
Text-prefix sidecar & Model-independent learned observer \\
Final-layer latent controller & Frozen-LM hidden-state baseline \\
Multi-layer latent controller & Parallel depth-aggregated observer \\
Voluntary + watchdog hybrid & Tests independent rescue of voluntary misses \\
Self-RAG-style main model & Prior-art main-model comparator, faithfully labeled \\
Oracle timing/type & Upper bound before evidence and R2 errors \\
\bottomrule
\end{tabularx}
\end{table}

If exact Self-RAG training is infeasible, call the condition Self-RAG-inspired.
Main-model training is confined to that declared comparator. Text and latent
sidecars keep the generator frozen. Match controller complexity and runtime cost
where possible; activation access alone earns no native-interface claim.

\section{Independent Watchdog and Multi-Layer Control}
Voluntary control may share the generator's blind spots. Compare voluntary-only,
watchdog-only, and hybrid operation. Report error overlap and
\[
\operatorname{WRR}=\frac{\text{voluntary misses caught safely by watchdog}}
{\text{voluntary misses requiring intervention}}.
\]
Watchdog Rescue Rate matters only alongside false activation, unnecessary
retrieval, abstention cost, and lead time.

Reuse Paper 3's layer sweep and aggregation ladder: final layer, best
development-selected layer, fixed logit averaging, voting, learned scalar
weights, and a tiny MLP over logits. START and END may use different frozen
mixtures. Plot decodability by depth and test whether \textsc{retrieve} emerges
earlier than \textsc{verify}; probe accuracy alone is descriptive, not causal
benefit.

\section{R2 Resolution and Runtime Enforcement}
R2 compares the existing Paper 2.5 heuristic, token-aware BM25, capability or
source-native rules, broadcast, and oracle source. Learned source routing
remains gated on a real heuristic-to-oracle gap. Registry/backend changes must
not retrain R1.

Evidence statuses include supported, contradicted, unverified, stale,
ambiguous, conflict, and unavailable. When required evidence is not authorized
or sufficient, policy may force abstention, qualification, conflict reporting,
or further retrieval. Measure evidence override when generation contradicts
authoritative active state. An optional evidence-integration LoRA may be trained
only after trigger quality is established and must remain factorized from the
watchdog.

\section{Metrics and Cost Frontier}
Report START/END/type metrics and FAR; Early and Late Catch rates;
missed-before-commitment; lead tokens; WRR and error correlation; unsupported
commitment rate (UCR); Authorized Commitment Coverage; unnecessary retrieval;
evidence override; abstention precision/recall; source/R2 correctness; calls,
fanout, latency, and cost; calibration; and OOD/cross-model transfer. For
controller threshold $\tau$, report frontiers under
\[
C(\tau)=\lambda_{fp}C_{fp}+\lambda_{fn}C_{fn}
+\lambda_{lat}C_{lat}+\lambda_{tok}C_{tok},
\]
not F1 alone. Mandatory slices include high-confidence/high-risk and
low-confidence/low-risk trajectories, cases with evidence already available,
conflicts, and each epistemic role.

\section{Gates, Falsification, and Scope}
The source-specific learned router is \texttt{no\_go} while Paper 2.5's
heuristic closes the oracle gap. The generic latent watchdog is an exploratory
protocol because it tests timing and independent rescue, not source identity.
It advances only if it improves the UCR--cost frontier over semantic rules and
voluntary generic control with acceptable FAR and transfer.

Latent control is not justified if semantic rules match the frontier,
voluntary intent catches nearly all needs, watchdog errors correlate with
generator misses, false retrieval/abstention costs dominate rescues, or hidden
states do not transfer. Main-model, text-sidecar, latent, hybrid, and
task-dependent winners are all valid outcomes.

There is no transactional/backtracking contribution, learned source identity in
R1, silent promotion of uncertain interpretation, or PRA/native-KV dependency.
PRA may optionally search active task state before remote retrieval while
leaving unrelated old records unmaterialized. Paper 4 proceeds only if retained
Paper 3/3.5 state makes provenance, dependency, retraction, or transactional
scope failures concrete.

\section{Execution Order}
Build the claim-level timing freeze from Paper 1.5 trajectories; reuse Paper 3
control code; run semantic and FLARE baselines; train the final-layer probe;
sweep and aggregate layers; compare voluntary, watchdog, and hybrid control;
add active-evidence sufficiency and enforcement; test cross-model/OOD transfer;
and only then consider integration LoRA. Until those artifacts exist, this is a
registered plan rather than evidence of an epistemic watchdog advantage.

\begin{thebibliography}{9}
\bibitem{flare} Z. Jiang et al. Active Retrieval Augmented Generation.
\emph{EMNLP}, 2023. \url{https://arxiv.org/abs/2305.06983}.
\bibitem{selfrag} A. Asai et al. Self-RAG: Learning to Retrieve, Generate, and
Critique through Self-Reflection. \emph{ICLR}, 2024.
\url{https://arxiv.org/abs/2310.11511}.
\end{thebibliography}
\end{document}
